\begin{abstract}
本文研究一类带开放指标的张量网络收缩（TNC）问题：目标不是查询单一标量元素，而是生成完整的输出张量。在误差评估、可视化与跨系统接口等场景中，下游真正需要的是完整的稠密输出张量。此时，问题的关键在于如何在保持输出边界不变的前提下，降低生成该输出所需的内部收缩代价。为此，本文提出自适应草图（Adaptive Sketching）方法用于近似和加速TNC：先对输出索引空间做分块，再在交互图上构造划分树，并沿树递归维护分隔状态；方法始终保持输出边界显式存在，只对内部接口做自适应的低秩近似，并在内部合并处通过基于算子访问的随机化草图压缩避免显式形成膨胀接口。

实验主要围绕多种经典张量网络展开，关注三个问题：是否存在可用的精度--时间折中、收益是否确实来自收缩阶段、以及近似强度能否稳定控制误差。结果表明，在代表性实验中，我们提出的方法将相对误差控制在 $4.81\times 10^{-3}$ 以内，并相对直接TNC获得 $13.98\times$ 到 $84.05\times$ 的加速；相比之下，非自适应地采用固定低秩虽然更快，但相对误差达到 $0.65$--$0.98$。实际实验也表明，加速几乎完全来自收缩阶段；机制分析进一步说明草图与输出分块是主要效率来源。实验进一步显示，近似强度可以控制误差和时间的权衡。
\end{abstract}
